\section{Qemu Memory Management}
\label{chap:QEMU_Memory}

Understanding QEMU's memory architecture is essential for the implementation 
of new devices. At the highest level, QEMU operates as a user-space application on 
the host operating system. This means that QEMU is responsible for modeling the 
entire memory hierarchy of a guest machine, including RAM, memory-mapped I/O 
devices (MMIO), ROM, and specialized controllers, all within a single process on 
the host operating system. This architecture enables transparent emulation of 
unmodified guest operating systems and firmware by presenting a faithful abstraction 
of physical hardware memory layout, serving these three critical functions:

\begin{enumerate}
    \item \textbf{Abstraction:}     
        \begin{itemize}
            \item Despite its internal complexity, this memory architecture presents a 
            unified interface for various types of memory.
        \end{itemize}
    \item \textbf{Flexibility:} 
        \begin{itemize}
            \item Enable dynamic memory management, region reallocation, and support 
            for specialized memory drivers.
        \end{itemize}
    \item \textbf{Performance:} 
        \begin{itemize}
            \item Maintaining translation caches and utilizing efficient lookup structures to 
            minimize address translation overhead.
        \end{itemize}
\end{enumerate}


To achieve what is described in the previous paragraph, QEMU's memory 
management system operates through two main layers of abstraction: the \textbf{MemoryRegions} 
and the \textbf{AddressSpaces}. While each serves a distinct function in the address 
translation process, understanding them is essential to understanding how guest 
memory accesses are transformed into host operations and how device emulation is 
achieved.

Unfortunately, while comprehensive for developers who want to implement board 
or SoC models, QEMU's official documentation doesn't provide a complete 
understanding of the memory architecture. It is assumed that the board/SoC developer 
audience doesn't need to know the details of how the memory subsystem uses the 
address space and memory region tree created by the board and SoC code. This 
design choice might be pragmatic, as board/SoC developers rarely need to modify 
the memory subsystem; they only need to use its APIs. Therefore, the documentation 
only offers a general overview of QEMU's memory regions and address spaces 
and provides much more detailed information on how to instantiate and connect 
memory regions, create device register blocks, connect them to address spaces, and 
define read/write controllers.

Under the previous context. In QEMU, the Memory is modelled as an acyclic 
graph of MemoryRegion objects. Sinks (leaves) are RAM and MMIO regions, while 
other nodes represent buses, memory controllers, and memory regions that have 
been rerouted. An AddressSpace is a mapping from guest addresses to memory 
regions—it tells QEMU which MemoryRegion handles each address in the guest's 
address space. There are multiple types of memory regions, all represented by a 
single C type MemoryRegion\cite{Qemu_MMIO_Ops}:

\begin{itemize}
    \item \textbf{RAM Regions:}
        \begin{itemize}
            \item Represent ordinary guest RAM. These region is simply a range of host memory that can be made available to the guest. 
        \end{itemize}

    \item \textbf{MMIO Regions:}
        \begin{itemize}
            \item Represent memory-mapped I/O devices accessed via callback functions. Each read or write causes a callback to be called on the host.  
        \end{itemize}

    \item \textbf{ROM Regions:}
        \begin{itemize}
            \item Represent read-only memory (firmware, bootloaders). Reading from ROM accesses host memory directly, but writes either are silently ignored or invoke callbacks (in ROM device variant).
        \end{itemize}

    \item \textbf{IOMMU Region:} 
        \begin{itemize}
            \item An IOMMU region translates addresses of accesses made to it and forwards them to some other target memory region. As the name suggests, these are only needed for modelling an IOMMU, not for simple devices. 
        \end{itemize}

    \item \textbf{Container Regions:}
        \begin{itemize}
            \item Represent hierarchical groupings of other memory regions, organizing complex address spaces. Useful for grouping several regions into one unit. For example, a PCI BAR may be composed of a RAM region and an MMIO region.
        \end{itemize}

    \item \textbf{Alias Region:} 
        \begin{itemize}
            \item Allow a subsection of one MemoryRegion to be mapped at a different location in the address space, enabling memory remapping without duplication. Aliases reference an offset and size within their target region.
        \end{itemize}

    \item \textbf{Reservation region:} 
        \begin{itemize}
            \item Represent a reservation region is primarily for debugging. It claims I/O space that is not supposed to be handled by QEMU itself. The typical use is to track parts of the address space which will be handled by the host kernel when KVM is enabled.
        \end{itemize}
\end{itemize}
